% NAL 2338, batch 5 — Gallica views 81–100.
% Pass: Opus 5.5 (claude-opus-5-5), 2026-09-23 — first pass, unchecked against the views by a human.
%
% Read from the local mirror only (archives/nal-2338/f0081–f0100.jpg), as
% thumbnails, four full-width bands per written view at 2400 px, half-width
% bands for the fast hand of v. 89–90, and tight crops for doubtful words.
% Greyscale and colour microfilm-style scans, portrait, about 4316 × 5740,
% one page per view, often with the edge of the next leaf at one side.
%
% The catalogue gives no dating: \dating{} is left empty, and this pass
% infers none. The dates written on the leaves (1636, 1643, 1644, 1645, 1647)
% are transcribed where they stand and are not a dating of the volume.
%
% Continuity. v. 81 does NOT open a piece: it begins mid-word (« miam,
% Gnomonicam ») inside a Latin letter begun on an earlier view (batch 4).
% v. 100 does NOT close one: it stops mid-sentence (« altius ; ») and the
% text runs on at v. 101 (batch 6).
%
% What the twenty views hold, piece by piece. The leaves name no writer
% except in the titles of v. 83 and v. 97, which name Roberval; nothing is
% attributed here beyond what the leaves carry.
%
%   v. 81–82   End of a Latin letter to a « vir clarissime », in a posed
%              regular copying hand (hand 1): the writer's Mechanics in eight
%              books (« octo contignationibus »), a prodromus printed in 1636
%              « gallico idiomate », a criticism of the correspondent's
%              « prop. 1a Lib. pi de motu grauium descendentium », paradoxes
%              of the press and of ropes, the centre of percussion of a
%              circular sector, and the « leges » of a correspondence. Stops
%              mid-sentence at v. 82 (« verum existimasse »), the lower third
%              of the page blank; the end is not in this batch.
%   v. 83–84   « Epistola Ægid. Personerii(?) de Roberual ad R. P.
%              Mersennum » (hand 2, upright, looped capitals, short t): a
%              critique of Torricelli's propositions, numbered in the margin
%              10–14 (the cochlea, centres of gravity of parabolic solids —
%              ratios 11:5, 13:7, 15:9, 17:11 —, the cycloid « nostra
%              Trochoides », with a struck passage on Beaugrand; solids of
%              conic sections; the acute hyperbolic solid, art. 14).
%              A pencil note, modern: « Il faut commencer par cette lettre ».
%   v. 85      A leaf with the figure of art. 14 only (hyperbola, asymptotes
%              BA, BC, points lettered then numbered 3–28). Described in a note.
%   v. 86      Blank verso, the figure of v. 85 showing through. Not
%              transcribed.
%   v. 87–88   Art. 14 continued, same hand 2: inscribed and circumscribed
%              cylinders, a corollary, and « Secundo », the infinite
%              hyperbolic solid equal to a cylinder, by exhaustion. Stops
%              mid-sentence at the foot of v. 88 (« alter cylindrus »); the end
%              is not in this batch.
%   v. 89–90   « Methodus Torricelli de dimensione infinitarum parabolarum
%              … » (hand 3, a very fast, heavily abbreviated Latin cursive,
%              the hardest of the batch: 310 of the 315 \ill{}): exponents of
%              « applicatae » and « diametrales », « Lemma 11 », « linea
%              Robervaliana », quadrature of parabolas (v. 89); « Lemma I »,
%              « Lemma 2 », « Theorema », then the spiral « quam spiralem …
%              seu Torricellianam appellabimus », its tangent and its length
%              equal to a parabola's (v. 90). Only the lettering, numbers and
%              technical terms are sure; much is \ill{}. The order of v. 89 and
%              v. 90 is not settled (v. 89 cites « 2m lemma », which v. 90
%              states).
%   v. 91      A leaf with a figure only, lettered for the « Theorema » of
%              v. 90 (staircase of rectangles under a curve). Note.
%   v. 92      Blank verso (a faint trace of the figure). Not transcribed.
%   v. 93–94   « Extraict d'une lettre du Sieur Torricelli », marginal « Al
%              S.r Ricci », then « Extraict d'une autre lettre du mesme au
%              mesme sur le mesme suject » (hand 4, French titles, Italian
%              text): the mercury tubes A, B, the weight of the air
%              (« 500 miglia d'aria »), water to 18 braccia; the reply to the
%              first objection, the cylinder of wool ABCD with the lead E.
%              Two figures of tubes (v. 93), one of the wool cylinder (v. 94).
%   v. 95      Recto of a letter-cover, blank (the address shows through).
%              Not transcribed.
%   v. 96      Its verso: the address « Au Reuerend Pere / Le Reuerend Pere
%              Mersenne Religieux ez(?) minimes de la Place Royale », upside
%              down, in a hand not met elsewhere in the batch; a modern
%              shelfmark « Bibl. Nat. n. a. l. 2338 ». Which letter it covered
%              is not shown.
%   v. 97–100  « De Vacuo Narratio Æ.i P.i de Roberual Ad Nob. virum D. des
%              Noyers » (hand 5, a posed fair copy): Valerianus Magnus's
%              booklet of July 1647, Torricelli's letter to Ricci of late
%              1643, Mersenne's journey, the experiments of « D. de Paschal »
%              at Rouen in January–February « huius anni » with tubes of 40
%              feet, the 2 7/24 feet of mercury, the « lagena » of 18 pounds,
%              the writer's own experiment at Paris with a drop of air, the
%              Aristotelians and the « spiritus » of mercury. Runs on past
%              v. 100.
%
% Views not transcribed: 86, 92 (blank versos, show-through checked at
% contrast); 95 (blank recto of the cover). v. 85 and 91 carry figures only
% and are given as a \page with a descriptive \note.
%
% Numbers on the leaves.
%   - One ink series, top right of the rectos only, one number per leaf across
%     five hands and across leaves that carry only a figure or nothing:
%     39 (v. 81), 40 (v. 83), 41 (v. 85), 42 (v. 87), 43 (v. 89), 44 (v. 91),
%     45 (v. 93), 46 (v. 95), 47 (v. 97, with a second figure struck beside
%     it), 48 (v. 99). Versos carry none (a reversed « 39 », « 42 » showing
%     through at v. 82, 88). It crosses hands and blanks and behaves as the
%     foliation; but it is ink, and, as for NAF 5176, Fr. 12357 and NAF 5161,
%     no \folio{} is written until a human has checked it against the
%     library's certificate (not in this batch).
%   - The marginal numbers 10.–14. on v. 83–84 number the articles of the
%     letter (Torricelli's propositions), not leaves.
%   - A large « Y » alone at the top left of v. 83: an order or bundle mark
%     of unknown meaning.
%   - In the figures, numbers 3–28 (v. 85) and 6-8, 8-21 … (v. 87–88) name
%     points, not quantities; « 25 » on v. 88 names a magnitude.
%
% Conventions. The long s is transcribed s. The copyists' u and v are kept
% as written (« uero », « vt », « vn »), with the abbreviations: « \&ca »,
% « oẽs », « nrãm », « ipsũ », « coĩs », « V. S.\textsuperscript{a} ». Mixed
% fractions are set $2\frac{7}{24}$. The sign « \& » of hand 4 is set \&.
% No sign of the period needed a substitute glyph.

% Coordinator's reconciliation (2026-09-23), after all six batches:
%   One ink series of leaf numbers runs top right of the rectos across every
%   hand and piece: 1–10 (v. 2–20), 11–20 (v. 22–40), 21–29 (v. 42–59),
%   30–38 with « 34 bis » for the figure slip (v. 61–79), 39–48 (v. 81–99),
%   49–50 (v. 101–103). It ends on « 50 » at the last written leaf, as the
%   certificate of v. 1 has it (« Volume de 50 Feuillets / 1er Octobre
%   1888 »): it is the library's count. Being ink, it gets no \folio{}, in any
%   batch. The bold series on the letter of v. 24–35 (1–11) is that letter's
%   own pagination (batch 2).
%   Seams: v. 60 ends « vniuersa- » and v. 61 opens « lius »; batch 3 had set
%   the catchword into v. 60, and now ends where the page does. v. 100 runs
%   into v. 101 (« altius ; » / « in aliis autem humilius »), the « De Vacuo
%   Narratio » (v. 97–103). The « non » sign is set « nl » in batch 1 (hand B,
%   an n with a barred ascender) and « n^o » in batch 2 (the fast cursive, an
%   n with a raised loop): two hands, two signs, left as each pass read them.

\documentclass[11pt,a4paper]{article}
% The preamble shared by every transcription of the mathematicians' manuscripts in Gallica.
%
% It rests on one idea: **uncertainty compounds, it does not resolve.** A
% transcription that smooths over a doubtful word has destroyed the only thing
% it was made for — a reader must be able to tell, at every point, what was on
% the page and what was supplied. The apparatus macros below are therefore the
% critical apparatus, not decoration.
%
% The same commands are recognised by `scripts/render.mjs`, which produces the
% site's reading view, and by `scripts/tei.mjs`. A macro added here must be
% added there too, or rendering fails loudly — which is the wanted behaviour.
%
% Comments here are English, like the rest of the repository. The documents
% this preamble sets are French, and that is deliberate: see the skills.

\ProvidesPackage{germain}[2026/09/15 Transcriptions of mathematicians' manuscripts in Gallica]

\usepackage{amsmath,amssymb,amsthm}
\usepackage{mathrsfs}
% Kept from the parent preamble so the renderer's diagram support stays
% available; these manuscripts draw no commutative diagrams, and nothing here uses it.
\usepackage{tikz-cd}
\usepackage[english,french]{babel}
\usepackage{fontspec}

% --- Greek ------------------------------------------------------------------
% The text face stays fontspec's default, Latin Modern, which has no Greek:
% XeTeX drops every glyph it lacks with a line in the log and nothing on the
% page, and the Greek words the manuscripts quote vanished from the PDFs. So
% the Greek and Greek Extended blocks get a XeTeX character class of their own,
% and entering or leaving a run of it switches the family to CMU Serif — the
% same Computer Modern design, with polytonic Greek — and back. Only the
% family changes: a Greek word in \textit or \textbf keeps its shape and series.
% The transitions to and from every other class carry the tokens that class
% has towards an ordinary letter, so French spacing before « ; » or inside
% « » still holds after a Greek word. No grouping, so no brace or \endgroup
% can come out unbalanced; maths is left alone, since XeTeX inserts nothing
% there. Kept identical in `exercices.sty`.
\makeatletter
\newfontfamily\germainGreekfont{cmun}[
  Extension=.otf, NFSSFamily=germaingreek,
  UprightFont=*rm, ItalicFont=*ti, SlantedFont=*sl,
  BoldFont=*bx, BoldItalicFont=*bi, BoldSlantedFont=*bl]
\newXeTeXintercharclass\germain@greek
\def\germain@greekrange#1#2{%
  \@tempcnta=#1\relax
  \loop
    \XeTeXcharclass\@tempcnta=\germain@greek
    \ifnum\@tempcnta<#2\advance\@tempcnta\@ne\repeat}
\germain@greekrange{"0370}{"03FF}
\germain@greekrange{"1F00}{"1FFF}
\def\germain@greekprev{\rmdefault}
\def\germain@greekin{%
  \edef\germain@greekprev{\f@family}\fontfamily{germaingreek}\selectfont}
\def\germain@greekout{\fontfamily{\germain@greekprev}\selectfont}
\def\germain@greekjoin{%
  \XeTeXinterchartokenstate=\@ne
  \@tempcnta=\z@
  \loop
    \ifnum\@tempcnta=\germain@greek\else
      \XeTeXinterchartoks\germain@greek\@tempcnta=\expandafter{%
        \expandafter\germain@greekout\the\XeTeXinterchartoks\z@\@tempcnta}%
      \XeTeXinterchartoks\@tempcnta\germain@greek=\expandafter{%
        \the\XeTeXinterchartoks\@tempcnta\z@\germain@greekin}%
    \fi
    \ifnum\@tempcnta<4095\advance\@tempcnta\@ne\repeat}
% babel-french sets its own classes, and saves and restores the interchar
% state, each time a language is selected: join again after it does.
\addto\extrasfrench{\germain@greekjoin}
\addto\extrasenglish{\germain@greekjoin}
\AtBeginDocument{\germain@greekjoin}
\makeatother

\usepackage{geometry}
\geometry{a4paper,margin=25mm,marginparwidth=28mm}
\usepackage{marginnote}
\usepackage{xcolor}
\usepackage[normalem]{ulem}
\setlength{\emergencystretch}{3em}
\usepackage{hyperref}
\hypersetup{colorlinks=true,linkcolor=black,urlcolor=black,pdfborder={0 0 0}}

% --- Batch metadata ---------------------------------------------------------
% Taken as they stand by the rendering script, which shows them at the head of
% the reading view. They fix what the file claims to cover. The macro names are
% the parent project's, so that the scripts stay shared; \folder here means the
% volume's slug — `fr-9115` — and \pages counts Gallica views.

\newcommand{\folder}[1]{\gdef\grothFolder{#1}}
\newcommand{\batch}[1]{\gdef\grothBatch{#1}}
\newcommand{\pages}[2]{\gdef\grothFirst{#1}\gdef\grothLast{#2}}
\newcommand{\foldertitle}[1]{\gdef\grothFoldertitle{#1}}

% The holder's own shelfmark, printed as they write it: « Français 9115 ».
\newcommand{\shelfmark}[1]{\gdef\grothShelfmark{#1}}

% The dating, copied verbatim from the holder's catalogue — never inferred
% here. « XIXe siècle » is the BnF's; a pass that dates a leaf from its content
% says so in a \note{}, as its own inference.
\newcommand{\dating}[1]{\gdef\grothDating{#1}}

% The status notice, printed in the title block and repeated as a diagonal
% watermark on every page of the PDF. The manuscripts are in the public
% domain; what this line declares is not a right but a quality — an
% unverified first pass by a machine. The skills write the line; a file
% without it gets no watermark, so leaving it out is visible in the output.
\newcommand{\watermark}[1]{\gdef\grothWatermark{#1}}

\gdef\grothFolder{?}\gdef\grothBatch{}\gdef\grothFirst{?}\gdef\grothLast{?}%
\gdef\grothFoldertitle{}\gdef\grothShelfmark{}\gdef\grothDating{}\gdef\grothWatermark{}

% --- The résumé -------------------------------------------------------------
\newenvironment{resume}
  {\small\setlength{\parskip}{0.45em}}
  {\par\medskip\hrule\medskip}

% --- Keywords ---------------------------------------------------------------
% In English, closing the modernised reading's résumé: the modern vocabulary
% under which to search for what the leaves construct. The manifest extracts
% them to tag the volume — this line is the single source of the tags.
\newcommand{\keywords}[1]{\par\smallskip\noindent
  {\footnotesize\textsc{Keywords} --- \textit{#1}}\par}

% --- The critical apparatus -------------------------------------------------

% The Gallica view number, in the margin. This is the anchor that lets a
% disagreement over a formula be settled by looking at *one* image rather than
% a volume of 750. The site uses it to turn the facsimile as the transcription
% is scrolled. A view is one scanned image: usually one side of a leaf.
\newcommand{\page}[1]{%
  \marginnote{\footnotesize\color{gray!70}\textsf{v.\,#1}}%
  \label{page:#1}\ignorespaces}

% The BnF's pencilled foliation, where the leaf carries it and a pass can read
% it: « 198r ». This is what the literature cites — « ff. 198r–208v » — and
% the one line that turns a citation into a view. Recorded, never inferred:
% a folio a pass cannot read is not written.
\newcommand{\folio}[1]{%
  \marginnote{\footnotesize\color{gray!50}\textsf{f.\,#1}}\ignorespaces}

% The modernised reading's coarser counterpart to \page{}: which views a
% section is drawn from.
\newcommand{\pagerange}[2]{%
  \marginnote{\footnotesize\color{gray!70}\textsf{v.\,#1--#2}}%
  \ignorespaces}

% Illegible: never guessed.
\newcommand{\ill}{\textcolor{red!60!black}{[\,\ldots\,]}}

% A doubtful reading: the word is given, and flagged as uncertain.
\newcommand{\uncertain}[1]{\uline{#1}}

% An editorial addition — a missing letter, an implied word.
\newcommand{\add}[1]{\textcolor{blue!50!black}{[#1]}}

% Struck out by the writer. What was crossed out is part of the document.
\newcommand{\struck}[1]{\sout{#1}}

% The transcriber's note, on the state of the page or the mathematics.
\newcommand{\note}[1]{\par\smallskip{\small\color{gray!60!black}\textit{[#1]}}\par\smallskip}

% A marginal note of hers, distinct from the body of the text.
\newcommand{\marginal}[1]{\par\smallskip\noindent\hspace{1em}%
  {\small\color{gray!60!black}#1}\par\smallskip}

% --- Title ------------------------------------------------------------------
% The title block is French, like the documents themselves.

\AddToHook{shipout/background}{%
  \ifx\grothWatermark\empty\else
    \begin{tikzpicture}[remember picture, overlay]
      \node[rotate=52, scale=3.4, align=center, text=gray!18,
            font=\sffamily\bfseries]
        at (current page.center) {\grothWatermark};
    \end{tikzpicture}%
  \fi}

\AtBeginDocument{%
  \begin{center}
    {\large\bfseries Manuscrits de mathématiciens (Gallica) ---
      \ifx\grothShelfmark\empty\grothFolder\else\grothShelfmark\fi}\\[2pt]
    {\itshape\grothFoldertitle}\\[4pt]
    {\small vues \grothFirst--\grothLast
      \ifx\grothBatch\empty\else\ (lot \grothBatch)\fi}\\[2pt]
    \ifx\grothDating\empty\else
      {\small\color{gray!60!black}Datation du catalogue : \grothDating}\\[2pt]
    \fi
    \ifx\grothWatermark\empty\else
      {\small\bfseries\color{red!45!gray}\grothWatermark}\\[2pt]
    \fi
    {\footnotesize\color{gray!60!black}Transcription établie d'après les images IIIF de Gallica.
     Source gallica.bnf.fr / Bibliothèque nationale de France.\\
     Les lectures incertaines sont soulignées, les illisibles notées [\,\ldots\,],
     les ajouts entre crochets ; \textsf{v.} numérote les vues de Gallica,
     \textsf{f.} la foliotation au crayon de la BnF.}
  \end{center}
  \vspace{1em}}

\endinput


\folder{nal-2338}
\shelfmark{NAL 2338}
\batch{5}
\pages{81}{100}
\dating{}
\watermark{Lecture automatique\\non vérifiée}
\foldertitle{Papiers divers de Mersenne.}

\begin{document}

\page{81}
\note{Recto d'une copie latine au net, d'une écriture posée et régulière, qui commence au milieu d'un mot : le texte continue une page qui précède cette vue. En haut à droite, à l'encre, « 39 » (série qui court d'une main à l'autre sur les rectos de ce lot ; voir l'en-tête). Le feuillet est monté ; le texte s'arrête aux trois quarts de la page, au milieu d'une phrase, qui se poursuit au verso (vue 82).}

miam, Gnomonicam, et Geographiam.

Plura quidem feci, quàm quæ comprehendere dictis

In promptu mihi sit.
\note{ces deux lignes sont écrites en retrait, comme une citation en vers.}

Sed illa omnia vulgaria æstimo. Attamen, dic quibus in terris Luna minori spatio quàm 24 horarum nostrarum communium, bis oriatur, aut bis occidat, ejusdem horizontis respectu : facile quidem theorema, sed quod primâ fronte impossibile multis videatur. At Mechanicam à fundamentis ad fastigium, nouam extruximus ; rejectis omnibus, præter paucos admodum, antiquis lapidibus, quibus illa constabat ; ita vt nunc octo contignationibus, hoc est totidem libris, absoluatur. Primus est de centro virtutis potentiarum in vniuersum ; an detur tale centrum, et quibus potentijs conueniat, quibus verò minimè. Secundus, de libra ; vbi de æquiponderantibus. Tertius, de centro virtutis potentiarum in specie. Quartus, de fune mira continet. Quintus, de instrumentis et machinis. Sextus, de potentijs quæ in diuersis medijs agunt. Septimus, de motibus compositis. Octauus denique, de centro percussionis potentiarum mobilium. In his omnibus nulla admitto noua postulata, sed tantùm ea quæ vulgò recepta sunt apud authores : quod sanè exequi, quàm non facile opus sit, testes sunt quotquot huc vsq; de grauibus super planis inclinatis existentibus egerunt : inter quos et ipse haberis vir clarissime, qui prop. 1\textsuperscript{a} Lib. p\textsuperscript{i} de motu grauium descendentium, ad id demonstrandum, nouo postulato vsus es, quod quiuis non facilè concesserit, quia pondera quæ proponis, non librâ rigidâ et rectâ, vt fieri solet, sed fune molli ac perfectè plicabili inuicem aligantur. Nos autem ad hoc, librâ vtimur modo vsitato dispositâ, cujus beneficio propositionem illam non aliter demonstramus, quàm aut vectem, aut axem in peritrochio : eam autem jam ante quindecim annos inuenimus, atque anno 1636 tanquam Mechanicæ nostræ prodromum, prelo commisimus atque vulgauimus, sed gallico idiomate. Neque etiam eum tantùm casum considerauimus qui solus ab omnibus attenditur ; cùm scilicet potentia pondus in plano inclinato positum retinens, agit per lineam directionis ipsi plano parallelam ; sed et dum eadem linea directionis aliam quamcumq; positionem obtinuerit : quo pacto, ratio ponderis ad potentiam infinitè mutatur. Ibi autem quiddam demonstrauimus, quod multis omninò paradoxum visum est : nempe, si intelligatur prelum aliquod duobus planis parallelis perfectè rigidis constans, quod ita disponatur vt ejus plana horizonti non sint parallela : tunc, quantâcunque potentiâ prematur prelum illud, planis semper perfectè planis, ac parallelis inter se remanentibus ; illa nullum pondus inter se retinebunt ; sed illud pondus propriâ grauitate statim labetur inter ipsa plana, atque idem à prelo sese liberabit,
\note{la lettre initiale de « Gnomonicam » est le grand G bouclé de cette main, le même que dans « Geographiam ». Le « v » de « vectem » est repassé sur une autre lettre, peut-être un r. Le trait qui traverse « ante » paraît la barre prolongée du t, non une rature. « quamcumq; », « vsq; » : abréviation de -que, laissée telle quelle.}

\page{82}
\note{Verso du feuillet de la vue 81, même main ; la phrase laissée en suspens au bas du recto continue ici. En haut à gauche, un « 39 » inversé et pâle est le chiffre du recto vu par transparence. Cachet rond de la bibliothèque (« R.F. ») à droite, vers le bas de la page.}

nisi aliunde retineatur. Hæc quidem ad quintum nostrum librum pertinent. Libet autem ex quarto quoque hæc addere. Si tres potentiæ totidem funibus ad communem nodum religatis agentes, (nodus est quoduis punctum in fune) æquilibrium constituant : tunc describi poterit triangulum cujus centrum grauitatis sit nodus ipse, tres autem anguli ad tria funium puncta alicubi terminentur. (infinita quidem describerentur triangula, sed omnia similia) Erunt autem tunc tres potentiæ in eadem ratione cum tribus rectis à centro trianguli ad tres angulos terminatis ; ita vt qualibet potentia homologa sit ei rectæ quæ in fune ipsius existit. Si quatuor potentiæ non existentes in eodem plano, totidem funibus ad communem nodum religatis agentes, æquilibrium constituant : tunc quod supra de triangulo dictum est, de quadam pyramide tetragona verum erit. Hinc aliud paradoxum. funis horizonti minimè perpendicularis quantâ vi tendatur, si perfectè plicabilis, nullo modo autem rigidus ex se existat ; imposito quocumque vel minimo pondere, aut si ipse ex se grauis esse intelligatur, flectetur necessariò, vel rumpetur ; nec viribus vllis fieri poterit vt rectus euadat. Similiter, tres, vel quotcunque funes ad communem nodum religati, totidem potentijs in eodem plano existentibus, quod planum horizonti non sit perpendiculare, quibuscunque viribus tendantur ; imposito quocunque vel minimo pondere, vel si ipsi funes per se graues esse intelligantur, nunquam tamen poterunt eò adduci vt in eodem plano consistant. Tandem etiam, ex octauo libro illud habebis. Omnis sectoris circuli semicirculo non \struck{\uncertain{maj}} majoris, circa centrum circuli circumuoluti, existente axe motus ad planum ejusdem circuli siue sectoris, perpendiculari, centrum percussionis siue impetus, in recta angulum sectoris bifariam diuidente quæsitum, sic reperietur. Vt chorda arcus sectoris ad ipsũ arcum, ita tres quadrantes semidiametri circuli ad rectam inter ipsius circuli centrum et centrum percussionis sectoris interceptam. Ex tali centro, quod extra sectorem aliquando existet, si impetus sectoris eo modo moti quo dictum est, excipiatur, productâ ad id rectâ angulum bifariam diuidente, si centrum illud extra sectorem excurrerit ; erit impetus ille maximus omnium qui ex quouis puncto in eadem recta existente excipi possunt.
\note{« paradoxum » est suivi d'un point, et « funis » commence la ligne suivante sans majuscule nette. Le mot biffé avant « majoris » est un premier « maj » abandonné, sous une rature épaisse. « ipsũ » : ipsum, avec le tilde de la nasale.}

De his et alijs agemus in posterum, si ita tibi placuerit vir clarissime, postquam, litibus valere jussis, solidam inierimus amicitiam, quam, vt spero, non recusabis : illius autem leges, quod ad literarum commercium attinet, tales sunto. Nihil tentandi gratiâ scribam. Quicquid scripsero, nisi de eo dubitare me, aut illud quærere scripsero, verum existimasse
\note{La page s'arrête ici, au milieu de la phrase, le dernier tiers de la page restant blanc. La vue suivante (83) ouvre une autre pièce ; la suite de cette phrase n'est pas dans ce lot.}

\page{83}
\note{Début d'une nouvelle pièce, sur un feuillet dont le bord supérieur est déchiré. Une autre main que celle des vues 81--82 : une cursive droite, à grandes capitales bouclées, qui remplit les fins de ligne d'un long trait. En haut à gauche, un grand « Y » isolé, marque d'ordre ou de liasse dont le sens n'est pas établi. Sous le titre, à droite, à l'encre, « 40 » (la série des vues 81, 83, 85 \ldots{} ; voir l'en-tête). À droite, dans un cadre, au crayon et d'une main moderne : « Il faut commencer par cette lettre », note de classement qui n'est pas du texte. Dans la marge de gauche, les numéros « 10. » et « 11. » en regard des articles.}

Epistola Ægid. \uncertain{Personerii} de Roberual ad R. P. Mersennum.
\note{le nom se lit « Persone- » suivi de deux jambages qui peuvent être n ou r (cette main fait l'r comme un v ou un n), puis d'un i surmonté de deux points.}

Reuerende Pater

Ex propositionibus Clarissimi Toricellj eas tantum examinandas censui, quas non nisi ab Egregio Geometra profectas esse iudicabam. quapropter prætergressis octo primis circa Sphæram, et solida eidem inscripta et circumscripta, quarum examen quemuis uel mediocriter uersatum, fugere non posse existimaui : nonam aggressus sum quæ est de dimensione cochleæ, quam, ut ardua est, ita ueram esse certissimâ demonstratione perspexi, ita ut ex ea unica Authorem inter præstantes huius sæculi Mathematicos annumerare non uerear. quodque fortassis mirere nihil referi magis-ne an minus inter se distent spiræ ipsius cochleæ, modo idem sit semper triangulum a quo describatur. Sed et etiam si ipsum triangulum moueatur tantum ad motum \uncertain{parallegrogrãmi} non autem motu progressiuo, ita ut idem triangulum absolutâ conuersione in seipsum redeat ; eodem modo se res habebit, nec mutabitur propositio.
\note{cette main fait le t court, presque comme un i : « ut », « et », « est » se lisent souvent « ui », « ei », « esi » ; on les a donnés ut, et, est. « referi » est écrit ainsi (pour refert). Le mot lu « parallegrogrãmi », avec le tilde de la nasale sur le second a, est écrit ainsi, pour parallelogrammi.}

\marginal{10.} De centro grauitatis parabolæ inueniendo a priorj, nulla supposita eius quadratura ; si ipse sic proponit, ut se inuenisse intelligat laudamus : si uero a nobis quærit dabitur illj non solum in parabola conica, quam quadratiuam appellamus, quia in ea quadrata ordinatim applicatarum inter se sunt ut portiones diametri : sed etiam in parabola cubica, in quadrato-quadratiua \&ca. atque in earum solidis ; siue ipsæ parabolæ circa suos axes, siue circa tangentes ad extremitatem axis, siue circa aliquam ex ordinatis ad axem conuertantur, et geniti inde solidi siue fusi parabolici dimidium plano ad ipsius axem erecto resectum proponatur : et multa alia de quibus, si aliquando res postulabit, fusius agemus. Nunc uero hoc indicasse sufficiat. In dimidio fuso parabolico quadratiuo centrum grauitatis axem diuidere in duas portiones quarum ea quæ ad uerticem ad eam quæ ad basim se habet \uncertain{ut} 11. ad 5. in cubico, ut 13. ad 7. In quadrato-quadratiuo, ut 15. ad 9, in quadrato-cubico, ut 17. ad 11. atque ita in infinitum addendo semper 2, ad singulos præcedentis rationis terminos : prætereo rationes solidorum ipsorum ad cylindros quibus inscribuntur, quas oẽs inuenimus et quarum speculatio forsan minime spernenda uiro Clarissimo uidebitur.
\note{avant « 11. ad 5. », le mot a la forme d'un « b2 » : on le donne « ut », comme dans les trois rapports qui suivent, où il s'écrit « ui ». « cylindros » s'ouvre sur une ligature à longue queue, comme un z, que cette main répète à la vue 87. « oẽs » : omnes.}

\marginal{11.} In Cycloide Toricellj agnosco nrãm Trochoidem, nec recte percipio
\note{la phrase continue au verso (vue 84). « nrãm » : nostram.}

\page{84}
\note{Verso du feuillet de la vue 83, même main. Le bord supérieur est déchiré ; aucun chiffre en haut de la page. Dans la marge de gauche, les numéros « 12. », « 13. », « 14. » en regard des articles.}

quomodo ipsa ad Italos peruenerit nobis nescientibus \struck{nisi fortan a \uncertain{I.} de Beaugrand missa fuerit, cui \ill{} solenne erat aliorum inuenta mutatis vocabulis et suppressis Authorum nominibus propalare}. quod si illa tanto uiro (\struck{Toricellum intelligo}) placuerit, lætor. spero autem breui fore ut eadem in lucem emittatur cum suis tangentibus, cumque solido ex conuersione illius circa basim genito, forsan et circa axem : neque id tantum in prima trochoide cuius basis æqualis esse ponitur circumferentiæ rotæ genitricis ; sed etiam in quauis alia trochoide siue prolata, siue contracta ; atque in socijs earumdem.
\note{trois lignes et demie sont biffées d'un trait continu, de « nisi fortan » à « propalare » ; la lecture en est sûre sauf le mot après « cui », commençant par h, sous une rature plus lourde, et l'initiale du prénom de Beaugrand, une capitale bouclée qui peut être I ou J. « fortan » est écrit ainsi (pour forsan). « Toricellum intelligo », dans la parenthèse, est barré d'un trait.}

\marginal{12.} Propositio de solido à qualibet sectione conj circa axem circumuoluta descripto, atque ad conum eidem inscriptum unica enunciatione collato, elegantissima est et uerissima ; demonstrauimus enim. Nec ei inferior est ea quæ sub eadem figura habetur de centro grauitatis ipsorum solidorum, quam etiam demonstrauimus : quod si ambas duabus tantum demonstrationibus ostenderet nihil uideo quod in hac materia desiderarj possit : sed uereor ne positis Authorum demonstrationibus ipse inde propositiones suas deduxerit : quod etiamsi ita esset, tamen non parum laudis mereretur : neque enim cuilibet \struck{cong} contingit, aliorum inuentis addere tanti ponderis propositiones.
\note{après « demonstrauimus enim » vient une sorte de parenthèse ouvrante isolée, sans parenthèse fermante.}

\marginal{13.} Eiusdem fere argumenti est sequens propositio de frusto sphærico duobus planis parallelis secto, de quo nihil dicimus, quia in eo non immorati sumus.

\marginal{14.} Omnium elegantissima est decimaquarta, cuius demonstrationem hic addere libet cuperemque ualde scire utrum in idem cum clarissimo uiro medium inciderim, uel diuersum. Igitur in figura cuius constructionem ex ipsius Toricellj propositione notam esse suppono existente B centro hyperb. asymptotis BA, BC, ad angulos rectos solido autem quouis DEFG, terminato, ut propositum est ; primum ostendamus tale solidum medium proportionale esse inter duos cylindros eiusdem altitudinis cum solido puta recta AH, quorum unius basis sit circulus DE, alterius uero FG ; ex hac enim cætera demonstrabuntur. Inter BA et BH media proportionalis sit BT ; tum inter BA, et \struck{media} BT, media quoque proport. sit
\note{la figure de ces lettres est seule sur le feuillet suivant (vue 85). La phrase continue en haut de la vue 87.}

\page{85}
\note{Feuillet qui ne porte qu'une figure, sans texte, pleine page, à l'encre, en traits pleins et pointillés ; en haut à droite, « 41 » (la série des rectos ; voir l'en-tête). Il est relié entre le verso de la vue 84 et la suite de la lettre (vue 87), et ses lettres sont celles de la démonstration de l'article 14 (BA, BC, AH, FG \ldots). Un axe vertical AB. Deux courbes pointillées, symétriques de part et d'autre de l'axe, qui s'écartent en descendant, coupées par des horizontales dont les extrémités et le milieu sont lettrés, de haut en bas : D A E ; I K L ; M N O ; P Q R ; S T V ; x y z ; 13, 3, 14, 4, 16, 5, 15 ; 9, 6, 10, 7, 12, 8, 11 ; F H G. Au pied, une horizontale lettrée 23, B, 24, C, que l'axe touche en B ; deux pointillés obliques partent de B et montent vers « 19 » à gauche et « 20 » à droite ; des verticales pointillées descendent de 18 vers 23 et de 17 vers 24, d'autres joignent 13, 9, F, 21, 22. En haut à gauche de la partie basse, « 26. 28. 27. » ; à droite, « 25. ». Dans l'article 14 (vue 84), B est le centre de l'hyperbole et BA, BC ses asymptotes : les deux courbes paraissent une hyperbole et sa symétrique par rapport à l'asymptote BA, qui s'en rapprochent en montant ; que les nombres 3 à 28 nomment des points une fois les lettres épuisées est une lecture de la figure. Une tache grise à gauche de la figure.}

\page{87}
\note{Suite de l'article 14 de la lettre des vues 83--84, même main, après le feuillet de la figure (vue 85, dont le verso, vue 86, est blanc). En haut à droite, « 42 » (la série des rectos ; voir l'en-tête). Cette main écrit « cylindrus » avec une ligature initiale qui ressemble à un g ou un z. Les points de la figure au-delà de l'alphabet sont nommés par des nombres ; les tirets entre deux nombres (6-8, 8-21) sont de la main du texte et joignent les deux extrémités d'un diamètre ou d'un rectangle.}

BN ; atque inter BT, et BH esto B4 : Item inter BA et BN sit BK ; inter B et BT, sit Bq, inter BT, et B4, sit By ; inter B4, et BH, sit B7. atque ita tot continue inueniantur mediæ quot libuerit, sic enim erunt quoque continue proportionales differentiæ ipsarum H7, 74, 4y, \&ca. usque ad ultimam KA, et in eadem ratione primarum. patet autem hac raone eò deuenirj posse ut cylindrus cuius basis circulus FG, altitudo autem ultima differentia KA, minor sit quouis spatio solido dato : jam per puncta 7, 4, y, T, \&ca. ducantur plana ad rectam AB, erecta solidum secantia secundum circulos quorum diametrj 6-8, 3-5, xz, SV, \&ca. parallela ipsi FG : patet quoque ex natura hyperbolæ, proportionales esse rectas FH, 6-7, 3-4, xy, ST, et reliquas in eadem ratione, sed inuersa primarum BH, B7, B4, \&ca. Denique inscribantur et circumscribantur ipsi solido totidem cylindri quot sunt differentiæ H7, 7-4, 4y, \&ca. Sintque inscripti 8-21, 5-10, z 14, \&ca. ; circumscripti uero F 11, 6-15, 3-17, \&ca. constat ergo omnes circumscriptos simul \marginal{+ Superare omnes inscriptos simul,} minori spatio quam cylindro altitudinis KA, et basis FG, hoc est minori spatio quouis proposito. præterea cylindrus basis SV, et altitudinis AH, est medius proportionalis inter cylindros eiusdem altitudinis, sed basium DE FG. diuidatur ipse medius in cylindros eiusdem basis SV, sed altitudinum H7, 7-4, 4y, yT, \&ca. usque ad ultimam altitudinis AK, qui ultimus major quidem est primo inscripto 8-21, sed minor circumscripto F, 11, quod sic ostendimus. quoniam recta ST, media proport. est inter DA, FH, major erit ratio circuli medij SV, ad circulum 6-8, quam recta DA, ad rectam 6-7 : at idem circulus medius SV, ad circulum FG, minorem habebit rationem quam eadem recta DA, ad eamdem 6-7, ut autem DA, ad 6-7, ita H7, ad AK ; ergo circulus medius SV, ad basim quidem inscripti 6-8, majorem habet rationem, ad basim uero circumscripti FG, minorem quam altitudo \uncertain{coĩs} inscripti, et circumscripti H7, ad altitudinem ultimi medij AK ; eodem modo demonstrabimus
\note{« inter B et BT » : la lettre après le premier B manque, un blanc la remplace (d'après la suite des moyennes, on attendrait BN ; c'est une inférence). « Bq » : la lettre est un q minuscule ; la figure a un Q capital sur l'axe entre N et T. « raone » est écrit ainsi (ratione). La marginale, écrite à gauche, est appelée par une croix placée après « simul » ; elle complète la phrase : les circonscrits surpassent les inscrits d'un espace moindre que le cylindre de hauteur KA. Devant « circumscripto F, 11 » la plume a d'abord formé un s. « coĩs » : communis, abréviation lue avec doute. La phrase continue au verso (vue 88).}

\page{88}
\note{Verso du feuillet de la vue 87, même main, bord supérieur déchiré ; en haut à gauche, le « 42 » du recto par transparence. La démonstration de l'article 14 continue.}

cylindrum altitudinis NK basis uero circuli medij SV majorem quidem esse secundo inscripto 5-10 ; minorem uero secundo circumscripto 6-15. atque ita de reliquis ordine sumptis. patet igitur tandem totum cylindrum medium omnibus quidem inscriptis simul sumptis majorem esse, omnibus uero circumscriptis minorem. cætera persequi apud uos inutile fuerit.

Corollarium

Patet autem manifestò positis rectis BH, B7, B4, By, \&ca continue proport. et factâ constructione eadem, diuidi totum solidum hyperb. FG, ED, in portiones continue proportionales in eadem quidem sed inuersa ratione rectarum ipsarum BH, B7, B4, \&ca. quæ portiones erunt FG, 8-6, 6-8-5-3, 3-5-ZX \&ca. quia qui ipsis portionibus æquales erunt cylindri proportionales erunt in ratione propositâ. quæ proprietas eximia est.

Secundo intelligamus solidum hyperbolicum BA, uertex A, infinite productum esse, atque idem secarj quouis plano 3-5, ad rectam BA erecto in puncto 4, ac circulum constituente cuius diameter 3-5 tum super hac base, circulo 3-5, esto cylindrus 3-5-24-23, cuius altitudo sit B4. Dico talem cylindrum æqualem esse solido hyperbolico super basi 3-5, constituto atque infinite uersus A extenso.

Alias uel cylindrus major est solido uel minor. esto primum major, si fierj potest, et excessus esto magnitudo 25, ita ut solidum hyperbol. una cum spatio 25, intelligatur æquale esse cylindro proposito 3-5-24-23. Jam intelligatur cylindrus quidam cuius altitudo B4, semidiameter uero basis Lq, ita ut hic cylindrus minor sit spatio 25, sit autem Lq perpend. ad BA, atque interjecta inter hyperbolam, et assymptoton hoc enim fierj potest. tum fiat ut B4, ad Bq, ita Bq, ad BA, et terminetur solidum hyperb. circulo DAE ; Erit ergo ex prædemonstratis, solidum 3-5-ED, æquale cylindro altitudinis A4, basis uero semidiametrj \uncertain{Lq}, addantur inæqualia solido quidem, spatium 25, cylindro uero alter cylindrus
\note{« 25 » nomme ici une grandeur (l'excès du cylindre sur le solide) ; sur la figure (vue 85), « 25. » est écrit seul, hors des lignes, en haut à droite, comme « 26. 28. 27. » à gauche. Dans « Lq », le L est la capitale bouclée de cette main et le q une minuscule ; la troisième occurrence, après « semidiametrj », a une initiale qui pourrait être un P. « hyperbolam » : la fin du mot est peu nette. La page s'arrête ici au milieu de la phrase ; la vue 89 ouvre une autre pièce, d'une autre main, et la suite de la démonstration n'est pas dans ce lot.}

\page{89}
\note{Nouvelle pièce, sur un feuillet d'un autre papier, d'une troisième main : une cursive latine très rapide et très abrégée, à peine formée, la plus difficile du lot ; seules les lettres des figures, les nombres et les termes techniques se lisent sûrement, et bien des mots ne sont donnés qu'avec doute. En haut à droite, « 43 » (la série des rectos ; voir l'en-tête). Dans la marge de gauche, quatre figures à la plume, en regard des paragraphes qu'elles servent : (1) un triangle de sommet C et de base AD, coupé par une courbe et par les horizontales OP et BI, avec L sur le côté droit ; (2) un rectangle lettré A, B, C, D, E, F, H, Z (?) où se croisent deux courbes ; (3) un parallélogramme E D C A, avec F et B, une courbe de E à B et les diagonales AD, AB ; (4) une figure lettrée E, H, P, A, O, L, I, F, B, G, C, avec une courbe de A à C par F et une droite de A vers E. Cachet rond de la bibliothèque au milieu de la page. La moitié inférieure est blanche.}

Methodus Torricelli de dimensione infinitarum parabolarum \uncertain{anno} \uncertain{1646} \uncertain{a Roberuallio}
\note{titre écrit au-dessus du texte, de la même main. Le nombre après « parabolarum » se lit « c46 » ou « 1646 » ; le dernier mot commence par « R » et paraît un nom propre. Rien sur le feuillet ne dit qui l'a écrit.}

Esto parabola quæcumque ABC, ad cuius punctum A ducta sit tangens. Applicetur AD ad diametrum ED, fiatque ut exponens applicatarum, qui exempli gratia sit 5, ad exponentem diametralium, qui exempli gratia sit 3, ita recta ED ad DC, dico rectam AE \ill{} \ill{} parabolæ ad aliud punctum præter A. Nam si fieri potest conueniat in B, et applicata BI, fiat ut ED ad DC, ita EI ad IO, sitque \ill{} recta CL, cui æquidistet OP.
\note{« quæcumque » est écrit en abrégé (q̃ʒ), comme plus bas ; « exempli gratia », « qui » : abréviations lues d'après le sens. Les deux mots après « rectam AE » ne sont pas lus (le premier pourrait être une négation). Le mot avant « recta CL » commence par b.}

Erit \uncertain{ergo} dignitas BI, hoc \uncertain{est in nostro casu} cuboquadratum BI, æquale ducto dignitatis IC in CL, nempe ducto cubi IC in quadratum CL ob \ill{} \ill{}. At idem cuboquadratum BI æquale ducto cubi IO in quadratum OP, ut \ill{}, \ill{} ducti IC in CL, et IO in OP sunt æquales, quod est impossibile, ut \ill{} \ill{}.

\uncertain{Quod} \uncertain{vero} \ill{} \ill{} \struck{\ill{}} sic : Cuboquadratum AD ad cuboquadratum BI \uncertain{ratio est} \ill{} rationis AD ad BI, siue rationis DE ad EI, siue cubi DC ad IO, et \uncertain{quadrati} AD ad BI, siue quadrati DE ad EI, ut quadrati CI ad EO, \ill{} quadrati CL ad OP ; \uncertain{ergo} cuboquadratum AD ad BI est ut ductus cubi DC in quadratum CL, ad ductum cubi IO in quadratum OP. \uncertain{Sed} antecedentia \uncertain{sunt} æqualia ob \ill{} \ill{}, \uncertain{ergo} et \uncertain{consequentia} \ill{}, \ill{} \ill{} \ill{}.
\note{un mot biffé d'une tache d'encre à la troisième place de la première ligne. La forme « cuboquadratum » (5\textsuperscript{e} puissance) est écrite en abrégé, « Cuboq\textsuperscript{tu} ».}

Secundum \uncertain{uero} ostendimus sic. Recta IO ad OE, est ut exponens diametralium ad differentiam exponentium, nempe ut 3 ad 2, quare ductus cubi IO in quadratum OE erit maximus (\uncertain{quod} \uncertain{prædemonstrauimus} \ill{} \uncertain{geometrice} \uncertain{siue} \ill{} \uncertain{per} \uncertain{algebram}) \ill{} ductus IO in OE ad ductum IC in CE, ita ductus IO in OP ad IC in CL. Quare et IO in OP maximus erit. Hæc \ill{} \ill{} methodus quam \struck{\ill{}} \ill{} \uncertain{de tangentibus}.

Lemma 11

Sit parallelogrammum, puta AB, CD, cui inscripta sit parabola (duæ parabolæ sint eiusdem speciei) \ill{} \ill{} \ill{}. Ponimus \ill{} figuras \ill{} \ill{} ut \ill{}, \ill{} \ill{} \ill{}, addita 3\textsuperscript{a} parabola in parallelogrammo BC, et \ill{} \ill{} duabus \ill{} \ill{} E, F, G, HP.

Ad quadraturam parabolarum \ill{} \ill{}, Esto parabola quæcumque ABC, cuius diameter BC, basis AC, parallelogrammum circumscriptum sit FC. Esto tangens AD, \ill{} in parallelogrammo FBDE alia parabola eiusdem speciei, siue \ill{} linea Robervaliana (\uncertain{nomine} \ill{} linea \ill{} \ill{} Robert. \ill{}) \ill{} \ill{} \ill{} aliqua parabolarum \ill{} \ill{} parabola eiusdem speciei, \ill{} \ill{} \ill{} \ill{}, \ill{} \ill{} \ill{} \ill{}) \ill{} parabola ABC æqualem trilineo ABE \ill{} \ill{} \ill{} \ill{} \ill{} solis \ill{} \ill{} \ill{} \ill{}.
\note{« Robervaliana » se lit sûrement, avec une capitale R repassée, au début de la sixième ligne du paragraphe ; « Robert. », abrégé, plus loin sur la ligne.}

\ill{} ut exponens applicatarum ad exponentem diametralium, ita recta DC ad CB ob tangentem \ill{} \ill{} lemma \ill{}, et ita \ill{} parallelogrammum EC ad FC, \ill{} ita \ill{} parabola EBD, ABC ad parabolam ABC ob 2\textsuperscript{m} lemma \ill{} \ill{} 19, 5\textsuperscript{i}, reliquum EBA siue parabola ABC \ill{} æquale erit ad reliquum FBA ut \ill{} \ill{} ad \ill{} \ill{} ut exponens applicatarum ad exponentem diametralium. \ill{} \ill{}, \ill{} \ill{} parallelogrammis \ill{} ad \ill{} parabolas \ill{} \ill{} \ill{} ad exponentem applicatarum.

\ill{} et alia \ill{} \ill{} \ill{} \ill{} quadratura omnium parabolarum, \ill{} \ill{} \ill{} \ill{} \ill{} \ill{}, \ill{} \ill{} \ill{} \ill{} \ill{} illa \ill{} linea, \ill{} \ill{} \ill{} \ill{} parabola \ill{} \ill{} \ill{}.

Esto parabolarum aliqua BAFC, cuius diameter AB, basis \uncertain{uero} BC, parallelogrammum inscriptum BL, linea Roberu. \ill{} \ill{} sit \uncertain{recta} APE \ill{} \ill{} quocumque puncto F in linea parabolica, sit applicata FI, tangens FH, et \ill{} \ill{} diameter æquidistet : \ill{} \ill{} DH applicatis æquidistans et \ill{} \ill{} ob definitionem \ill{}.

\ill{} DF recta ad FO erit ut HT ad BA, siue ob tangentem ut exponens applicatarum ad exponentem diametralium. \ill{} ita \ill{} \ill{} \ill{} simul \ill{} \ill{} recta DF, \ill{} \ill{} \ill{} \ill{}, \ill{} trilineum EDAFC, siue parabola ABC \ill{} æqualis, \ill{} ad \ill{} \ill{} \ill{}, \ill{} ad trilineum LAC, ut exponens applicatarum ad exponentem diametralium, quadratura \ill{} \ill{} \ill{}, \ill{} ut 3. Et \ill{} 4 alias \ill{} \ill{} Torricelli \ill{} tangentibus \ill{} per quadraturam, quæ \ill{} \ill{} \ill{} et \ill{} parabolarum et hyperbolarum, et plurimas aliorum \ill{}, \ill{} \ill{} \ill{}
\note{la page s'arrête sur ces mots, à mi-hauteur ; la suite est au verso (vue 90). Dans tout ce texte, les lettres des figures et les nombres sont sûrs ; les mots de liaison, réduits à des traits, ne l'étaient pas assez pour être donnés.}

\page{90}
\note{Verso du feuillet de la vue 89, même main très rapide et très abrégée, plus lisible cependant qu'au recto. Aucun chiffre en haut de la page. La page ne continue pas la dernière phrase du recto : elle s'ouvre sur « Lemma I » et « Lemma 2 », que le recto paraît invoquer (« ob 2\textsuperscript{m} lemma »), de sorte que l'ordre de lecture des deux pages n'est pas établi. Figures à la plume dans la marge de gauche : (1) un parallélogramme lettré D, B, E, A, I, F, C, O, avec une courbe et une diagonale ; (2) un rectangle lettré D, M, E, O, F, C, avec une courbe et des horizontales ; (3) un demi-cercle de diamètre BA, de centre C, avec les rayons CD, CF, CE ; (4) un triangle rectangle C, B, A, avec F sur l'hypoténuse, I, D, E, et un arc pointillé autour de A. Dans le texte, entre deux lignes, un petit croquis lettré B, C, A : un arc et un cercle pointillé autour de A. Au bas de la page, à droite, deux traits obliques en fin de texte.}

Lemma I

In figura Roberualli \ill{} \ill{} \ill{} \ill{} quodlibet parallelogrammum inscriptum AB \ill{} erit \uncertain{quam} FA \uncertain{circumscriptum} \ill{} figura \ill{}. Nam ducta tangente per A, similibusque figuris, æquale erit AB \ill{} AI, ergo \ill{} quod erat demonstrandum.
\note{« Roberualli » se lit sûrement, en abrégé ; la suite de la ligne n'est lue qu'aux lettres.}

Lemma 2

\uncertain{Circumscriptum} \uncertain{parallelogrammum} \struck{FE, quod \ill{} inscriptum \ill{} figuræ genitrici} CD maius erit quam aliud parallelogrammum FE, quod est inscriptum figuræ genitrici. Nam tangente ducta per O, et completa figurâ, erit CD æquale ipsi FM, ergo CD maius quam FE. quod erat demonstrandum.
\note{la ligne biffée répète, d'un trait continu, ce que la phrase dit ensuite ; un mot après « aliud », sous une rature, est suivi de « parallelogrammum » récrit.}

Theorema

Sumatur primo aliqua pars figuræ \ill{}, puta ABCDE, dico æqualem esse figuræ BGC. Aliàs maior \ill{} uel minor. Sit primo maior, secetur que BK sæpe bifariam, donec HE minor sit excessu. Tunc inscribatur in figura \uncertain{iuncta} BCE alia figura ex parallelogrammis æqualis RD, LM etc. quorum ultimum sit IO. Eritque inscripta figura, ob \ill{}, adhuc maior spatio BGC, quod est contra lemma \ill{}.

Sit deinde minor, et secetur BK sæpe bifariam, ut \ill{} sit HE minor defectu. Tunc inscribatur figuræ BCE alia figura constans ex parallelogrammis æque altis, et erit inscripta figura adhuc maior figura BCE, quod \ill{}. Nam ad figuram inscriptam trilineo BCE maior est quam alia quæuis figura inscripta ipsi BGC ob 2\textsuperscript{m} Lemma, \ill{} \ill{}, quod erat demonstrandum.

\ill{} \ill{} \ill{} (\ill{} \ill{} \ill{} \ill{}) dico figuram \ill{} æqualem esse. Aliàs aliqua \ill{} maior sit, sit maior APC, et \struck{\ill{} \ill{}} erit aliqua figura \ill{} pars, puta BGC æqualis alteri figuræ, \ill{} absurdum est.

Si uero parallelogrammum \uncertain{maius} ACEF, erit aliqua ipsius pars, puta BCE æqualis figuræ genitrici APC, absurdum est.
\note{ce paragraphe est souligné, et un trait le sépare du suivant.}

\ill{} \ill{} pro linea spirali nostra. Esto quæcumque \struck{figura} recta linea AB secta utcumque in C. Deinde erigatur perpendicularis CD, quæ sit media proportionalis inter BC, CA. Item recta CE sit media inter DC, CA, secetque bifariam angulum DCA. Amplius recta CF media inter DC, CE, secetque bifariam angulum DCE, et sic fiat semper ducendo medias proportionales quæ bifariam secent angulos. \ill{} omnia puncta A, E, F, D, B et quotcumque alia \ill{} per quæ transibit linea nostra, quam spiralem \ill{} seu Torricellianam appellabimus.
\note{« Torricellianam » est sûr ; le mot qui le précède, après « spiralem », n'a pas été lu.}

Aliter definitur. Si recta linea AB in \ill{} figura in plano æqualiter \ill{} uoluatur, et interim aliquod punctum C \ill{} linea \ill{} \ill{} mouendo \ill{} \ill{} \ill{} \ill{}, \ill{} \ill{} A, \ill{} spatia æqualibus temporibus peracta inter se proportionalia sint, tandemque \ill{} \ill{} \ill{} distantia \ill{} puncti à centro A, ita \ill{} spiralem \ill{}.

Peculiare \ill{} est huic lineæ, quod \ill{} perueniat ad centrum, \ill{} \uncertain{revolutiones} in \ill{} absoluere debebit. \ill{} longitudinem \ill{} \ill{} quorum \ill{}, si \ill{} \struck{\ill{}} \ill{} \ill{}. Nam si fuerit centrum A \ill{} radius AB, \ill{} BC tangens, et angulus ABC rectus, erit tangens BC æqualis \ill{} spirali BDA. Dato uero arcu quocumque BI, facillime ipsi recta linea æqualis abscindetur \ill{}.

Ad hoc \ill{} completa \ill{} methodus quæ \ill{} \ill{}. \ill{} \ill{}, \ill{} \ill{} \ill{} ad \ill{} \ill{} \ill{} \ill{} rectis. Spatium quoque \ill{}, \ill{} \ill{} \ill{} \ill{} \ill{} \ill{} \ill{} ad \ill{} \ill{} \ill{}. \ill{} \struck{\ill{}} spatia \ill{} \ill{} \ill{} \ill{} \ill{} \ill{} æqualia \ill{} spiralis Archimedis \ill{} \ill{} \struck{\ill{}} \ill{}, \ill{} \ill{} \ill{} spatio \ill{} \ill{} \ill{} \ill{} \ill{} \ill{} \ill{} \ill{} \ill{} \ill{} \ill{} \ill{}, \ill{} \ill{} \ill{} \ill{} spiralem \ill{} parabolas. \ill{} \ill{} \ill{} \ill{} \ill{} \ill{} \ill{} \struck{parabolas} hyperbolas et ellipses \ill{} \ill{} \uncertain{Mersenno} \ill{} \ill{} \ill{} de rectis abscissis à tangentibus \ill{} \ill{} \ill{}, \ill{} \ill{} lineam spiralem utriusque lineæ parabolicæ æqualem esse \uncertain{demonstrauimus}.
\note{les huit dernières lignes sont les plus rapides de la page. Les lettres « Archimedis » (lu « Archimedes » ou « Archimedis », au milieu de la quatrième ligne), « parabolas », « hyperbolas », « ellipses », « de rectis abscissis à tangentibus » et « lineam spiralem \ldots{} lineæ parabolicæ æqualem » sont lues ; au-dessus de « parabolas » biffé est écrit « hyperbolas ». Le mot lu avec doute « Mersenno » est à la fin de la ligne où se lisent « hyperbolas » et « ellipses », précédé d'une abréviation qui pourrait être celle de « Reuerendo Patre ». Le reste n'a pu être lu.}

\page{91}
\note{Feuillet sans texte, qui ne porte qu'une figure à la plume dans son tiers supérieur, le reste blanc ; en haut à droite, « 44 » (la série des rectos ; voir l'en-tête). À gauche de la vue, la fin des lignes de la vue 90, qui dépasse sous ce feuillet. Les lettres de la figure sont écrites couchées, comme si la figure devait être lue le feuillet tourné d'un quart de tour. Un grand rectangle K P A (le coin supérieur gauche marqué E, au-dessus F) ; une courbe descend de E vers C sur le côté inférieur, en passant par D, M, O ; une droite fine va de C vers A en passant par B ; entre la courbe et cette droite, une suite de rectangles en escalier, de plus en plus petits vers C, dont les sommets portent D, M, I, O ; à droite, des verticales descendent vers le côté KP, l'une jusqu'en G ; sur la ligne haute, B, B (ou 8), L, N, H. Près de E, une courte ligne de hachures. Les lettres B, C, E, G, H, K, D, M, I, O sont celles du « Theorema » de la vue 90 (BGC, BCE, BK, HE, RD, LM, IO) ; que cette figure lui serve est une lecture de la figure.}

\page{93}
\note{Nouvelle pièce, sur un feuillet isolé, d'une quatrième main : un titre en français, puis un texte en italien, d'une cursive régulière. En haut à droite, « 45 » (la série des rectos ; voir l'en-tête). Dans la marge de gauche, deux figures à la plume : (1) deux tubes de verre dressés dans une cuvette, l'un terminé en haut par une boule marquée E, l'autre droit ; l'horizontale AB les coupe à la même hauteur ; la cuvette est marquée D (le niveau de l'eau) et C, C (le mercure, hachuré au fond) ; les deux colonnes de mercure sont hachurées jusqu'à AB ; (2) un tube à boule, recourbé en U par le bas, lettré f, E, a (un petit rectangle accolé au tube), B et C. Le texte s'arrête aux trois quarts de la page, sur « \&c. », souligné d'un trait.}

Extraict d'une lettre du Sieur Torricelli.

\marginal{Al S.\textsuperscript{r} Ricci.}
\note{écrit dans la marge de gauche, en face du titre, de la même main.}

Noi habbiamo fatto molti vasi di vetro, \& anco canne come i segnati A, \& B, grossi o di collo lungo due braccia. questi pieni d'argento viuo, poi serratigli con vn dito la bocca \struck{\ill{}} \& riuoltati in vn vaso doue era l'argento viuo C, si vedeuano votarsi \& non succeder niente nel vaso che si votaua. Il collo però AD restaua sempre pieno fin all'altezza d'vn braccio \& $\frac{1}{4}$ \& vn dito di più. per mostrar poi che il vaso fosse perfettamente voto, si riempiua la catinella sotto posta d'acqua fino in D, \& alzando il vaso a poco a poco, si vedeua quando la bocca del vaso arriuaua all'acqua, descender quell'argento viuo del collo, \& riempirsi con impeto horribile d'acqua fin' al segno E affatto.
\note{« vn » devant « vaso » est écrit au-dessus de la ligne, avec un signe d'insertion ; la lettre effacée avant « \& riuoltati » est couverte d'une tache. Le premier « poco » est corrigé sur une autre lettre. La lettre E de « fin' al segno E » est écrite à la fin de la ligne, dans la marge de droite. Cette main écrit « et » par un petit signe (\&) et « un » « vn ».}

Il discorso si faceua mentre il vaso AE staua voto, \& l'argento viuo si sosteneua ben che grauissimo nel collo AC ; questa forza che regge quel argento viuo contro la sua naturalezza di ricader giù si è creduto fin' adesso che sia stata interna nel vaso AE, ò di vacuo, ò di quella robba sommamente rarefatta, ma io pretendo che ella si esterna, \& che la forza venga di fuori. Su la superficie del liquore che è nella catinella grauita l'altezza di 500 miglia d'aria. però qual marauiglia è se nel vaso CE, doue l'argento viuo non ha inclinatione, ne anco repugnanza, per non esserui nulla entro, \& vi sinalzi fin tanto che si equilibri con la grauità del aria esterna che lo \uncertain{costringe}. L'acqua poi in vn vaso simile ma molto più lungo salirà quasi fin' à 18 braccia. cio è tanto più del argento viuo, quanto l'argento viuo è più graue del acqua per equilibrarsi con la medesima cagione che spinge l'uno \& l'altro. Confermaua il discorso l'esperienza fatta nel medesimo tempo col vaso A, \& colla canna B, nei quali l'argento viuo si fermaua sempre nel medesimo horizonte AB, segno quasi certo che la virtù non era dentro, perche più forza \uncertain{hauerebbe} hauuto il vaso AE, doue era più robba rarefatta, \& attrahente, è molto più gagliarda per la rarefazion maggiore che quella del \uncertain{pochezzo} spatio B \&c.
\note{« si esterna » : le mot qui suit « ella » est bref et peut être un « sia » abrégé. « costringe » est lu avec doute (la plume a fait « cortinge »). Le mot lu « hauerebbe » est chargé d'une rature ou d'une surcharge. Le mot lu « pochezzo » reste douteux.}

\page{94}
\note{Verso du feuillet de la vue 93, même main : un second extrait, titre en français et texte en italien. Aucun chiffre en haut de la page. Au pied du texte, une figure à la plume : un cylindre ABCD (A et D en haut, B et C en bas) rempli de hachures ondulées (la laine), coupé par l'horizontale fg, surmonté d'un poids arrondi marqué E. Cachet rond « BIBLIOTHÈQUE NATIONALE / R.F. / MANUSCRITS » à droite de la figure.}

Extraict d'une autre lettre du mesme au mesme sur le mesme suject.

Quanto alla prima obiettione, io rispondo, se V. S.\textsuperscript{ria} quando \uncertain{indurrà} la lamina saldata che copra la superficie della catinella la \uncertain{indurrà} di maniera ch'ella tocchi l'argento viuo della catinella, io dico che quello inalzato nel collo del vaso resterà come prima solleuato non per il peso della esterna aria, mà perche quello della catinella non potrà dar luogo. Se poi V. S.\textsuperscript{ria} \uncertain{indurrà} si ch'ella pigli anco dentro dell'aria, io domando se quell'aria serrata dentro V. S.\textsuperscript{a} che sia nel medesimo grado di condensazione che la esterna. \& in questo caso l'argento viuo si sosterrà come primà (per l'esempio che darò adesso della lana) mà se l'aria che V. S.\textsuperscript{a} include sarà più rarefatta dell'esterna, all'hora il metallo solleuato descenderà alquanto. se poi fusse infinitamente rarefatta cioè vacuo, all'hora il metallo descenderebbe tutto pur che lo spatio serrato la potesse capire. Il vaso ABCD è vn cilindro pieno di lana ouero d'altra materia compressibile (diciamo anco d'aria) il qual vaso hà due fondi, BC stabile \& \struck{altro} AD mobile \& che sadatti, \& sia AD caricato sopra dal piombo E, che pesi $\frac{m}{10000000}$ libre. vedo che V. S.\textsuperscript{a} intenderà quanta violenza sia per sentire il fundo BC ; hora se noi spingeremo a forza vn piatto tagliente fg, si che entri \& tagli la lana compressa io dico che se la lana fBCg sarà compressa come prima, ancor ch'il fundo BC non senta più nulla del peso sopra posto del piombo E, in ogni modo patirà il medesimo che patiua prima, applichi V. S.\textsuperscript{a} che io non starò a tediarla più \&c.
\note{le verbe lu « indurrà » s'achève chaque fois sur un paraphe, et reste douteux. « V. S.\textsuperscript{ria} », « V. S.\textsuperscript{a} » : Vostra Signoria. Le poids du plomb est écrit en fraction, « m » au-dessus d'un trait, « 10000000 » au-dessous.}

\page{96}
\note{Verso du feuillet numéroté « 46 » (vue 95, dont le recto est blanc : l'adresse y transparaît). La feuille porte les plis d'une lettre fermée ; dans le panneau du haut, tête-bêche par rapport à la vue, une adresse en deux blocs, d'une main qui n'est aucune de celles des pièces voisines. Au bas de la page, à l'encre, une cote moderne « Bibl. Nat. n. a. l. 2338 », non transcrite. L'adresse ne dit pas de quelle lettre elle était le couvert.}

Au Reuerend Pere

Le Reuerend Pere Mersenne Religieux \uncertain{ez} minimes de la Place Royale
\note{« Reuerend » est écrit en abrégé les deux fois ; le mot lu « ez » est une ligature en forme de g.}

\page{97}
\note{Nouvelle pièce, sur un feuillet plus grand, d'une cinquième main : une copie latine au net, d'une écriture régulière et bien formée, qui remplit les fins de ligne d'un long trait. En haut à droite, « 47 » (la série des rectos ; voir l'en-tête), et à sa droite un autre chiffre effacé ou biffé, illisible. Cachet rond « BIBLIOTHÈQUE NATIONALE / R.F. / MANUSCRITS » au milieu de la page. Les millésimes sont soulignés par le copiste.}

De Vacuo

Narratio Æ.\textsuperscript{i} P.\textsuperscript{i} de Roberual Ad Nob. virum D. des Noyers.
\note{« Æ.\textsuperscript{i} P.\textsuperscript{i} » : deux initiales avec un i suscrit, abréviations des prénom et nom latins ; elles ne sont pas développées ici.}

De Vacuo quod in rerum natura facile dari permulti autumant : adducto ad id comprobandum, nobili experimento hydrargyri tubo inclusi, modo qui iam omnibus satis notus est : Ignoscat mihi R. P. Capucinus Valerianus Magnus, si dixero illum parum candide egisse in eo libello quem de hac re in lucem nuperrime emisit mense Julio huius anni 1647 dum celeberrimi huiusce experimenti ille primus author haberi voluit, quod certo constat jam ab anno 1643 in Italia vulgatum fuisse : ac ibidem, præcipue vero Romæ atque Florentiæ, celeberrimas inter eruditos de ea re viguisse controuersias, quas non potuit ignorare Valerianus qui circa eadem tempora illis in regionibus degebat, et cum doctis illis conuersabatur. Habeo ego Epistolam quam Clarissimus vir Euang. Toricellius magni Ducis Hetruriæ Mathematicus misit Romam ad amicum suum doctiss. virum Angelum Ricci sub finem anni 1643 Italice scriptam : quæ nihil aliud continet quam controuersiam inter duos illos viros egregios, qui, quod et fere omnibus accidit de tali experimento diuersa sentiebant. Ea autem Epistola cum quibusdam aliis, ab ipso Ricci missa est Parisios ad R. P. Mersenium ordinis Minimorum sub initium anni 1644 qui eo ipso anno Romam profectus est ubi et Valerianum ægrotantem tunc inuisit et cum eo quædam nostrorum opera recens edita communicauit. Sed et in eadem Epistola, vasa et Tubi quibus idem Toricellius vsus fuerat figuris exhibentur ; atque ex discursu apparet minime nouum tum fuisse illis experimentum, sed jam multoties repetitum. Tentatum quidem illud fuit ab ipso Mersennio statim post acceptam Toricellij de ea re Epistolam, verum defectu tubi ad id satis apti, nihil tunc fieri potuit ; ac non multo post ipse in Italiam profectus est : atque ob iter Florentiæ apud Toricellum, vasa et Tubos prædictos vidit et contrectauit. Idem autem reuersus sub finem anni 1645, rem omnem vulgauit : neque tamen eo anno aut sequenti tubos aptos Parisiis recuperare potuit ; tum quia ibi tales non fabricantur ; tum etiam quia ipse toto ferme eo tempore per Meridionales Regni Gallici partes peregrinatus est. Tandem ergo idem scripsit Rhotomagum ad amicos suos : ibi enim celeberrima habetur vitri et Crystalli officina. Sed antequam is inde Tubos haberet vulgatum fuerat et ibidem experimentum, ac plurimis modis, cum priuatim coram amicis, tum publice coram omnibus eruditis, multoties exhibitum a Nobiliss.\textsuperscript{o} viro D. de Paschal, mense Januario et Februario huius anni : neque id solum beneficio hydrargyri in Tubis minoribus, puta 3. 4. aut 5 pedum regiorum mensura nostra ; sed (quod mirandum multis videbatur) beneficio aquæ et vini in Tubis 40 pedum, ex Crystallo
\note{« Mersenium », puis « Mersennio » : les deux graphies sont sur la page. La phrase continue au verso (vue 98).}

\page{98}
\note{Verso du feuillet de la vue 97, même main ; la phrase laissée au bas du recto continue ici. Aucun chiffre en haut de la page. Le haut de la page, au-dessus de la première ligne, est blanc ; le titre du recto y transparaît.}

mira arte fabricatis, atque ad malum nauis alligatis, quod machinis ad id paratis ita libratum erat ut et attolli et deprimi ad usum requisitum facile posset. Et quidem occasio cur ad tantam altitudinem recurreret talis fuit.

Cum viderent omnes hydrargyrum merum, \struck{\ill{}} nulla aqua adhibita, nec in Tubo, nec in scutella subiecta ita in Tubo descendere, relicto superiori spatio Tubi veluti vacuo ut tamen semper occuparet in inferiori parte Tubi altitudinem pedum $2\frac{7}{24}$ proxime, mensuratam secundum perpendiculum supra superficiem hydrargyri in scutella subiecta contenti ; quantacunque pars Tubi in scutella immersa esset, et quantacunque extaret : illi, vt par est, in re noua, ignotis causis, in diuersas sententias abierunt. Quidam enim merum in tali spatio vacuum relinqui arbitrati, suam sententiam facile deffendebant contra aliorum obiectiones ; at nullatenus astruebant, deficientibus rationibus quibus nullum corpus, neque etiam subtilissimum, in tale spatium subingredi conuincerent. Quidam præcipue Peripatetici, qui Magistri sui Aristotelis, non rationibus quæ nullæ sunt sed uerbis mordicus adhærent ; contendebant paruulam quandam eamque forsan sensu imperceptibilem aeris guttulam remansisse, quæ postea eousque rarefieret ; ve sic naturæ laboranti opitularetur. At illi parum proficiebant : quia in diuersis Tubis latioribus et altioribus, eadem semper altitudo hydrargyri relinquebatur : Ideoque illa altitudo ad rarefactionem inducendam et conseruandam, æquales vires semper obtinebat ; cum tamen nunc minor, nunc maior foret rarefactio ; ad quam proinde, nunc minorem nunc majorem hydrargyri vim requiri consentaneum videbatur. Adde quod, adhibito tubo stricto quidem, sed qui in parte superiori dilataretur instar lagenæ quæ 18 libras hydrargyri contineret, eoque debito modo impleto atque disposito ut fit in aliis Tubis relinquebatur tota illa lagena ab hydrargyro vacua, cum parte superiore Tubi, donec illud in inferiore \struck{\ill{}} parte, sub altitudine prædicta pedum $2\frac{7}{24}$ quiesceret : nec quisquam in animum inducere poterat, posse insensibilem aeris guttulam a pondere 2 Unciarum hydrargyri in tubo remanentis, eam trahente, eò adduci ut in tantum lagenæ spatium rarefieret ; videbaturque, secundum illam Philosophiam, hydrargyrum debuisse potius altius ascendere, ad spatium illud occupandum : præcipue quia inclinato sensim tubo, donec altitudo verticis illius, perpendiculariter sumpta, esset pedum $2\frac{7}{24}$, aut minor, ille rursus totus hydrargyro replebatur : nihilque omnino sæpissime aut aliquando minimum quid extranei præter hydrargyrum, intra Tubum aut lagenam animaduertebatur
\note{« quantacunque pars », « quantacunque extaret » : écrits avec l'abréviation de -que. « altius » est écrit au-dessus de « ascendere ». Le mot biffé après « merum » est couvert d'une tache ; celui après « inferiore » est surchargé. « ve sic » est écrit ainsi (pour « ut sic » ?). La page s'arrête sur un long trait après « animaduertebatur », aux deux tiers de la hauteur, et le texte reprend en haut du feuillet suivant (vue 99) sans nouveau titre.}

\page{99}
\note{Recto d'un nouveau feuillet de la même copie (même main que les vues 97--98) ; le texte continue celui de la vue 98 sans titre. En haut à droite, « 48 » (la série des rectos ; voir l'en-tête). Le bas de la page est blanc.}

quod tamen extraneum, si quod erat, occluso Tubo, tum inuerso, et rursus recluso, ad os Tubi ascendebat, et infusa noua hydrargyri guttâ, fugabatur : Vt sic ; iterando aliquoties, si ita opus esset, merum hydrargyrum appareret ; quod tamen, erecto rursus Tubo supra scutellam, idem quod prius, saltem ad sensum, spatium vel \ill{} vacuum relinquebat : et iterum inclinato replebat. Cumque erecto Tubo, vacuum appareret ; si immisso in hydrargyrum scutellæ digito, Tubus occluderetur, et inde attolleretur ; apparebat semper vacuum, ut prius : at inuerso Tubo, cadebat hydrargyrum cum strepitu in fundum Tubi ; et statim vacuum versus digitum apparebat in superiori parte Tubi. Atque ita alternatim, inuerso Tubo ; ita ut digitus stricte illum occludens, nunc in superiori parte existeret ; nunc in inferiori : semper apparebat in parte superiori vacuum ; idque statim, nullâ interpositâ morâ : nec præter hydrargyrum quidquam intra Tubũ cernebatur. At antequam Tubus ex scutella extraheretur, si ille sic inclinaretur, ut apex illius ad altitudinem prædictam pedum $2\frac{7}{24}$ præcise perueniret nullum adhuc vacuum eo in statu cernebatur, cum tamen sic hydrargyrum totam suam altitudinem obtineret ; atque ideo et omnes suas vires exerceret, ad aerem si quis esset, trahendum ad se et sensibiliter rarefaciendum. Immo admissa sponte, in tali inclinationis statu, aeris guttâ (quod facile est) illa ab omnibus, et dum in Tubum per hydrargyrum ascenderet, et dum, eo superato, in summo ipsius nataret, facile cernebatur. Sique eadem satis ampla existeret, inclinato magis, ac magis Tubo, magis ac magis comprimebatur ; quippe ad ingenium redibat ; quia ab hydrargyro minus alto, minus trahebatur, minusque rarefiebat. at e contrario, dum Tubus sensim eleuaretur, fieretque hydrargyrum altius, atque ideo ad trahendum et rarefaciendum aerem, potentius ; tum guttâ aeris magis dilatabatur et tunc erecto ad perpendiculum Tubo hydrargyrum ad prædictam altitudinem, non omnino ascendebat. Quoque plus aeris prius admissum erat, dum Tubus inclinaretur, eo minus hydrargyri in Tubo erecto relinquebatur. Neque tamen, hac ratione, fieri vnquam potuit, ut aer magis rarefieret, quam ut octies amplius spatium occuparet, \struck{quod} quam illud quod eidem in libero statu deberetur. Quibus omnibus experimentis colligebatur guttam insensibilem a prædicto vi pedum $2\frac{7}{24}$ hydrargyri rarefactam, non potuisse totam lagenam 100000 majorem replere : neque etiam ab eadem vi minime mutata potuisse diuersos rarefactionis gradus, in eadem aeris guttam, ad diuersa spatia \struck{replenda} replenda, immitti : nec proinde talem opinionem stare posse, nisi apud eos qui, non luce veritatis sed tenebris ignorantiæ delectarentur ; vt andabatarum more, clausis oculis depugnarent.
\note{un mot après « vel » est couvert d'une tache. « minus » est écrit au-dessus de « trahebatur ». « quam » est écrit au-dessus de « quod » biffé. « replenda » est biffé en fin de ligne et récrit au début de la suivante. « Tubũ » : Tubum.}

\page{100}
\note{Verso du feuillet de la vue 99, même main ; nouveau paragraphe en haut de la page. Aucun chiffre en haut de la page. Le haut et le bas de la page sont blancs ; le titre d'un autre feuillet (« De Vacuo ») transparaît en haut.}

Præ cæteris autem, hoc mihi et multis aliis stupendum visum est, dum Parisiis circa tale experimentum exercerer. Inclinaui Tubum in scutella subiecta, sic ut ipse hydrargyro impleretur ; et per os quod inferius erat intra scutellam, quodque paululum eleuaui usque ad superficiem hydrargyri in scutella contenti, guttam aeris admisi, et statim idem os rursus in hydrargyrum immersi, ne plus aeris hauriret quam par esset. Ascendit ergo gutta aeris intra tubum, per medium hydrargyrum quousque in summo illius nataret : tum Tubum ad perpendiculum erexi ; et vacuum apparuit more solito hinc, supposito digito intra scutellæ hydrargyrum os Tubi occlusi et Tubum vna cum hydrargyro, aere et vacuo abstuli, omnia simul digito et manu sustinens atque ostentans iis qui astabant. hinc, quam potui celerrime tubum inuerti, ut digitus fieret superior, ac tunc velut in instanti apparuit vacuum versus meum digitum multo prius quam aeris gutta illuc ascenderet : hæc enim gutta videntibus omnibus, sensim per hydrargyrum ascendit, et ad vacuum superius peruenit, ut sic in inferiori parte Tubi merum relinqueretur hydrargyrum. Rursus autem celeriter Tubum inuerti et statim cadente hydrargyro in meum digitum non sine aliquo doloris sensu, propter impetum quamquam non altus esset casus, vacuum in superiori parte Tubi apparuit ; gutta autem aeris non nisi sensim eodem ascendit : et sic factum est toties quoties ego Tubum inuerti : Vnde omnes exclamarunt tale vacuum non esse aerem. Sed et idem ego cum gutta aquæ : et idem rursus cum gutta aquæ et gutta aeris expertus sum ; ac idem semper apparuit. Oportet autem ad hoc, Tubum non admittere plus quam 6 digitos vacui ; aliás hydrargirum vel digitum læderet, vel Tubum impetu frangeret ; quibus tamen periculis sublatis poterit eligi Tubus quicumque maior et elegantius euadet experimentum. Sed jam Rhotomagum redeamus.

Post Aristotelicos non pauci qui cæteris multo oculatiores \struck{nobis} sibi videbantur ; ex iis scilicet qui cogitata sua non minus quam nudam veritatem depereunt, ad hydrargyri naturam recurrentes, quam spiritibus abundare constat ; non cunctanter affirmarunt tale spatium in vertice Tubi relictum veluti vacuum, non re uera vacuum esse, sed totum a spiritibus illis obtineri. Verum illi primo quidem iisdem rationibus quibus Aristotelici, et quidem jure, refellebantur : quippe quod ex maiori aut minori copia hydrargyri, produci debuisset maior aut minor copia spirituum ; præcipue in diuersis Tubis diuersæ latitudinis et altitudinis : sicque hydrargirum in tubis quidem aliis, altius ;
\note{dans la marge de gauche, en face de « Post Aristotelicos », une grande croix d'un trait fin, peut-être d'une autre plume que le texte. « constat » : le t final est écrit au-dessus d'un e biffé. La page s'arrête au milieu de la phrase, sur « altius ; », aux deux tiers de la hauteur ; la suite est au-delà de ce lot (vue 101).}

\end{document}
